% Analytic matrix exponential of an infinitesimal planar rotation.
\clearpage
\section*{Verification case: matrix exponentials preserve the intended geometry}
\noindent A matrix exponential is not an elementwise exponential.
For the planar rotation generator
\[
\Omega=\begin{pmatrix}0&-\omega\\\omega&0\end{pmatrix},
\quad\text{with }\omega>0,
\]
powers of \(\Omega\) alternate between \(I\), \(\Omega\)
and their negatives because \(\Omega^2=-\omega^2I\).
Separating even and odd powers of the exponential series yields
\[
e^{t\Omega}
=I\cos(\omega t)+\frac{\Omega}{\omega}\sin(\omega t)
=\begin{pmatrix}\cos(\omega t)&-\sin(\omega t)\\
\sin(\omega t)&\cos(\omega t)\end{pmatrix}.
\]
Consequently, if \(x(0)=(1,0)^T\) and
\(t=\pi/(2\omega)\), then \(x(t)=(0,1)^T\).
The exact map is orthogonal and preserves squared lengths:
\(x(t)^Tx(t)=x(0)^Tx(0)=1\).

A single forward-Euler step with that full time increment instead
produces \(x_1=(1,\pi/2)^T\), whose squared norm is
\(1+\pi^2/4>1\). Both updates are derived from the same continuous
ODE \(\dot x=\Omega x\); they disagree because discretization has
changed the invariant. Smaller Euler steps reduce the error but do
not make a finite-step explicit Euler map exactly orthogonal.

\paragraph{Reconstruction.}
Implement the exponential of this \(2\times2\) generator using the
closed form and compare it with a trusted numerical matrix-exponential
library. Sweep \(t\) and measure orthogonality error
\(\lVert R^TR-I\rVert_F\). Compare Euler steps with a method known
to preserve the geometric structure under a stated model.
Do not infer that every matrix exponential is a rotation:
the generator's skew-symmetry is essential.
